\begin{frame}{변수 하나만 바꾼다: 8$\times$8 $\to$ 28$\times$28}
  \begin{itemize}
    \item \textbf{같은} digits 이미지를 8$\times$8에서
      \textbf{28$\times$28(784차원)으로 확대}해 다시 GBDT vs CNN
    \item 64차원: GBDT가 \textbf{0.006} 앞서던 것이
      $\Rightarrow$ 784차원: \textbf{CNN 0.961이 GBDT 0.944를
      0.017로 역전}
    \item 격차 \textbf{0.006 $\to$ $-$0.017} — 닫히고 심지어 뒤집힘
      = \kb{축 2(규모)}의 \textit{증거}
    \item 왜? 8$\times$8은 ``이미지''여도 64개 수치로 flattening되어
      \textbf{정형}이고, 28$\times$28(784개 픽셀)은
      \textbf{국소 구조를 공유하는 원시 픽셀}에 가까워지는 순간
      CNN의 \textbf{공유·평행불변 편향}이 일을 시작
  \end{itemize}
  \vskip0.5em
  \begin{block}{이 실험까지 세운 팀의 ``왜 DL'' 답변}
    ``64차원 정형에서는 GBDT 0.944가 상한이고 내 CNN(0.939)은 이를
    어깨너머로 확인했으며, 같은 데이터를 784차원으로 확대하자 CNN이
    오히려 0.961로 GBDT를 역전시켰다 — 내 문제가 차원·규모가 크다면
    (EMNIST, CIFAR) CNN의 우위는 이 추세의 연장선이라는 것을
    \textbf{숫자로} 보여준다.''
  \end{block}
\end{frame}
